\section{Discussion}
\label{sec:discussion}

\paragraph{Recovery is a vector, not a loss threshold.}
Perplexity remains the correct measure of the CPT objective, and the keep14
trajectory shows that continued optimization also produces a real MMLU gain.
What it does not provide is a certificate that a target behavior has returned.
The intact-base gap, interface audit, and closed-book results show why recovery
should be represented by a vector of measurements rather than by the direction
of a single loss curve.

\paragraph{Null baselines turn protocol sensitivity into evidence.}
Without Random, the 6.48-point keep14 content gain could be read as revealing
capability hidden by answer letters. The 11.28-point Random gain instead shows
that much of this interface is attainable as option-text compatibility after
training from scratch. This is a general evaluation principle: when a new
interface is introduced to reveal latent performance, its null floor is part
of the construct-validity test.

\paragraph{The ShortGPT gap is a design target.}
The 15.5-point letter-MMLU advantage of ShortGPT over keep14 is not merely a
warning about confounding; it identifies a productive next experiment. A
factorial study should independently vary inherited count, contiguous versus
BI selection, retention of the original final block, and fresh-tail insertion.
The much smaller 1.8-point content-score gap further suggests that these choices
may affect answer-interface calibration more strongly than option-text
likelihood. This hypothesis is directly testable with the existing paired
protocol.

\paragraph{A practical recovery audit.}
We recommend reporting six axes for post-intervention recovery:
(1) in-domain and, when possible, out-of-domain likelihood;
(2) target evaluations with sample counts;
(3) prompt/scoring/generation interface and a relevant null baseline;
(4) exact inherited/fresh structure and trainable modules;
(5) steps, token presentations, compute, and stopping rule; and
(6) training-run replication separately from item-level uncertainty.
These axes make recovery claims comparable even when pruning algorithms,
training budgets, and evaluation interfaces differ.
